% ============================================================
%  INFORME: MODELOS ARIMA Y SARIMA EN SERIES FINANCIERAS
%  Compilar con: pdflatex -shell-escape informe_arima_sarima.tex
%  (dos pasadas para el índice)
% ============================================================
\documentclass[12pt, a4paper]{article}

% ── Paquetes ──────────────────────────────────────────────
\usepackage[utf8]{inputenc}
\usepackage[T1]{fontenc}
\usepackage[spanish,es-nodecimaldot]{babel}
\usepackage{amsmath, amssymb, amsthm}
\usepackage{geometry}
\usepackage{graphicx}
\usepackage{float}
\usepackage{booktabs}
\usepackage{array}
\usepackage{xcolor}
\usepackage{listings}
\usepackage{caption}
\usepackage{subcaption}
\usepackage{hyperref}
\usepackage{fancyhdr}
\usepackage{titlesec}
\usepackage{tocloft}
\usepackage{enumitem}
\usepackage{mdframed}
\usepackage{setspace}
\usepackage{parskip}

% ── Geometría ─────────────────────────────────────────────
\geometry{
  top=2.5cm, bottom=2.5cm,
  left=3cm,  right=2.5cm,
  headheight=15pt
}

% ── Colores ───────────────────────────────────────────────
\definecolor{azulTec}{RGB}{0,84,166}
\definecolor{grisClaro}{RGB}{245,245,245}
\definecolor{grisOscuro}{RGB}{60,60,60}
\definecolor{verdeComentario}{RGB}{0,128,0}
\definecolor{rojoString}{RGB}{163,21,21}
\definecolor{azulKeyword}{RGB}{0,0,200}

% ── Estilo de código Python ───────────────────────────────
\lstdefinestyle{python}{
  language=Python,
  backgroundcolor=\color{grisClaro},
  basicstyle=\ttfamily\footnotesize,
  keywordstyle=\color{azulKeyword}\bfseries,
  stringstyle=\color{rojoString},
  commentstyle=\color{verdeComentario}\itshape,
  numberstyle=\tiny\color{grisOscuro},
  numbers=left,
  stepnumber=1,
  numbersep=8pt,
  showstringspaces=false,
  breaklines=true,
  breakatwhitespace=false,
  frame=single,
  rulecolor=\color{azulTec},
  tabsize=4,
  captionpos=b,
  xleftmargin=15pt,
  xrightmargin=5pt,
  aboveskip=10pt,
  belowskip=10pt
}
\lstset{style=python}

% ── Encabezado y pie de página ────────────────────────────
\pagestyle{fancy}
\fancyhf{}
\fancyhead[L]{\small\textcolor{azulTec}{\textbf{Modelos ARIMA y SARIMA}}}
\fancyhead[R]{\small\textcolor{grisOscuro}{Series de Tiempo Financieras}}
\fancyfoot[C]{\thepage}
\renewcommand{\headrulewidth}{0.4pt}
\renewcommand{\headrule}{\hbox to\headwidth{\color{azulTec}\leaders\hrule height \headrulewidth\hfill}}

% ── Formato de secciones ──────────────────────────────────
\titleformat{\section}
  {\Large\bfseries\color{azulTec}}{\thesection.}{0.5em}{}[\vspace{-0.3em}\color{azulTec}\rule{\linewidth}{0.6pt}]

\titleformat{\subsection}
  {\large\bfseries\color{grisOscuro}}{\thesubsection.}{0.5em}{}

\titleformat{\subsubsection}
  {\normalsize\bfseries}{\thesubsubsection.}{0.5em}{}

% ── Entornos de definición ────────────────────────────────
\newmdenv[
  backgroundcolor=grisClaro,
  linecolor=azulTec,
  linewidth=1.5pt,
  leftline=true, rightline=false, topline=false, bottomline=false,
  innerleftmargin=10pt, innerrightmargin=10pt,
  innertopmargin=8pt, innerbottommargin=8pt
]{defbox}

\newmdenv[
  backgroundcolor=white,
  linecolor=grisOscuro,
  linewidth=0.8pt,
  roundcorner=4pt,
  innerleftmargin=10pt, innerrightmargin=10pt,
  innertopmargin=8pt, innerbottommargin=8pt
]{notebox}

% ── Hiper-referencias ─────────────────────────────────────
\hypersetup{
  colorlinks=true,
  linkcolor=azulTec,
  urlcolor=azulTec,
  citecolor=azulTec,
  pdftitle={Informe: Modelos ARIMA y SARIMA},
  pdfauthor={Estudiante},
  pdfsubject={Series de Tiempo Financieras}
}

% ── Espaciado ─────────────────────────────────────────────
\setlength{\parskip}{6pt}
\onehalfspacing

% ══════════════════════════════════════════════════════════
\begin{document}
% ══════════════════════════════════════════════════════════

% ── PORTADA ───────────────────────────────────────────────
\begin{titlepage}
  \centering
  \vspace*{1.5cm}

  {\Huge\bfseries\color{azulTec} Modelos ARIMA y SARIMA\\[0.4em]}
  {\LARGE\bfseries Aplicados a Series de Tiempo Financieras\par}

  \vspace{1.2cm}
  \textcolor{azulTec}{\rule{0.7\linewidth}{1.5pt}}
  \vspace{1cm}

  {\large
    \begin{tabular}{rl}
      \textbf{Asignatura:}  & Series de Tiempo y Pronósticos \\[4pt]
      \textbf{Entrega:}     & Informe Final de Modelado       \\[4pt]
      \textbf{Fecha:}       & \today                          \\[4pt]
      \textbf{Herramienta:} & Python 3.x + LaTeX              \\
    \end{tabular}
  }

  \vspace{1.5cm}

  \begin{defbox}
    \centering\small
    Este informe presenta el marco teórico de los modelos AR, MA, ARMA, ARIMA y SARIMA,
    junto con tres ejemplos prácticos aplicados a datos financieros reales/sintéticos,
    incluyendo selección de parámetros mediante ACF, PACF, AIC y BIC,
    diagnóstico de residuos e interpretación de pronósticos.
  \end{defbox}

  \vfill
  {\small\textcolor{grisOscuro}{Documento generado con Python y \LaTeX}}
\end{titlepage}

% ── ÍNDICE ────────────────────────────────────────────────
\tableofcontents
\newpage

% ══════════════════════════════════════════════════════════
\section{Marco Teórico}
% ══════════════════════════════════════════════════════════

\subsection{Procesos Autorregresivos (AR)}

Un proceso \textbf{autorregresivo de orden $p$}, denotado AR($p$), expresa el
valor actual de la serie como combinación lineal de sus $p$ valores pasados
más un término de error:

\begin{defbox}
\begin{equation}
  y_t = c + \phi_1 y_{t-1} + \phi_2 y_{t-2} + \cdots + \phi_p y_{t-p} + \varepsilon_t,
  \qquad \varepsilon_t \sim \mathcal{N}(0,\sigma^2)
  \label{eq:ar}
\end{equation}
\end{defbox}

donde $c$ es una constante, $\phi_i$ son los coeficientes autorregresivos y
$\varepsilon_t$ es ruido blanco gaussiano. El proceso es \emph{estacionario}
si las raíces del polinomio característico $1 - \phi_1 B - \cdots - \phi_p B^p = 0$
se encuentran fuera del círculo unitario, siendo $B$ el operador de retardo
($B y_t = y_{t-1}$).

\subsection{Procesos de Media Móvil (MA)}

Un proceso \textbf{de media móvil de orden $q$}, denotado MA($q$), expresa la
serie como combinación lineal de $q$ errores pasados:

\begin{defbox}
\begin{equation}
  y_t = \mu + \varepsilon_t + \theta_1 \varepsilon_{t-1} + \theta_2 \varepsilon_{t-2}
        + \cdots + \theta_q \varepsilon_{t-q},
  \qquad \varepsilon_t \sim \mathcal{N}(0,\sigma^2)
  \label{eq:ma}
\end{equation}
\end{defbox}

Todo proceso MA($q$) es estacionario. La condición de \emph{invertibilidad}
exige que las raíces de $1 + \theta_1 B + \cdots + \theta_q B^q = 0$ estén
fuera del círculo unitario.

\subsection{Modelo ARMA($p$,$q$)}

El modelo \textbf{ARMA($p$,$q$)} combina los componentes AR y MA:

\begin{defbox}
\begin{equation}
  y_t = c + \sum_{i=1}^{p}\phi_i y_{t-i} + \varepsilon_t + \sum_{j=1}^{q}\theta_j \varepsilon_{t-j}
  \label{eq:arma}
\end{equation}
\end{defbox}

Usando el operador de retardo $B$, la ecuación~\eqref{eq:arma} se escribe
compactamente como:
\begin{equation}
  \Phi(B)\, y_t = c + \Theta(B)\, \varepsilon_t,
\end{equation}
donde $\Phi(B) = 1 - \phi_1 B - \cdots - \phi_p B^p$ y
$\Theta(B) = 1 + \theta_1 B + \cdots + \theta_q B^q$.

\subsection{Modelo ARIMA($p$,$d$,$q$)}

Muchas series financieras son \emph{no estacionarias}: su media o varianza
cambia en el tiempo. La solución habitual es aplicar $d$ diferenciaciones
hasta obtener estacionariedad. El modelo \textbf{ARIMA($p$,$d$,$q$)}
(\textit{AutoRegressive Integrated Moving Average}) resulta de aplicar un
modelo ARMA($p$,$q$) a la serie diferenciada $d$ veces:

\begin{defbox}
\begin{equation}
  \Phi(B)\,(1 - B)^d\, y_t = c + \Theta(B)\, \varepsilon_t
  \label{eq:arima}
\end{equation}
\end{defbox}

\begin{itemize}[leftmargin=*]
  \item $p$: orden autorregresivo (número de retardos AR).
  \item $d$: grado de integración (número de diferencias para lograr estacionariedad).
  \item $q$: orden de media móvil (número de retardos MA).
\end{itemize}

En la práctica, $d = 0$ para series ya estacionarias, $d = 1$ para series
con tendencia estocástica (como precios de activos) y $d = 2$ para series
con tendencia en la tendencia, aunque $d > 2$ es poco común.

\subsection{Modelo SARIMA($p$,$d$,$q$)($P$,$D$,$Q$)$_s$}

Cuando la serie exhibe \emph{patrones estacionales} de periodo $s$
(e.g., $s=12$ para datos mensuales), se añade un segundo nivel de diferenciación
y parámetros estacionales. El modelo \textbf{SARIMA} tiene la forma general:

\begin{defbox}
\begin{equation}
  \underbrace{\Phi(B)}_{\text{AR}}\;
  \underbrace{\Phi_S(B^s)}_{\text{AR estacional}}\;
  \underbrace{(1-B)^d}_{\text{Diferencia regular}}\;
  \underbrace{(1-B^s)^D}_{\text{Diferencia estacional}}\;
  y_t
  \;=\;
  \underbrace{\Theta(B)}_{\text{MA}}\;
  \underbrace{\Theta_S(B^s)}_{\text{MA estacional}}\;
  \varepsilon_t
  \label{eq:sarima}
\end{equation}
\end{defbox}

donde:
\begin{align*}
  \Phi_S(B^s) &= 1 - \Phi_1 B^s - \Phi_2 B^{2s} - \cdots - \Phi_P B^{Ps}, \\
  \Theta_S(B^s) &= 1 + \Theta_1 B^s + \Theta_2 B^{2s} + \cdots + \Theta_Q B^{Qs}.
\end{align*}

Los parámetros estacionales $(P,D,Q)_s$ tienen la misma interpretación que
sus análogos no estacionales, pero actúan sobre los retardos múltiplos de $s$.

\subsection{Funciones ACF y PACF}

\subsubsection{Función de Autocorrelación (ACF)}

La \textbf{autocorrelación} de orden $k$ mide la correlación lineal entre
$y_t$ e $y_{t-k}$:

\begin{equation}
  \rho_k = \frac{\gamma_k}{\gamma_0},
  \qquad \gamma_k = \text{Cov}(y_t, y_{t-k}) = E[(y_t - \mu)(y_{t-k} - \mu)]
  \label{eq:acf}
\end{equation}

La gráfica de $\rho_k$ versus $k$ se llama \emph{correlograma}. Las bandas
de confianza al 95\% se construyen como $\pm 1.96/\sqrt{n}$ (aproximación
para ruido blanco).

\begin{notebox}
\textbf{Interpretación para identificar modelos:}
\begin{itemize}[leftmargin=*,noitemsep]
  \item \textbf{AR($p$)}: ACF decrece geométricamente (cola larga); PACF corta en el rezago $p$.
  \item \textbf{MA($q$)}: ACF corta en el rezago $q$; PACF decrece geométricamente.
  \item \textbf{ARMA($p$,$q$)}: ambas funciones decrecen gradualmente.
  \item Si la ACF no decae (permanece alta), la serie no es estacionaria.
\end{itemize}
\end{notebox}

\subsubsection{Función de Autocorrelación Parcial (PACF)}

La \textbf{autocorrelación parcial} de orden $k$ mide la correlación entre
$y_t$ e $y_{t-k}$ \emph{eliminando} el efecto lineal de los valores
intermedios $y_{t-1}, \ldots, y_{t-k+1}$:

\begin{equation}
  \text{PACF}(k) = \text{Corr}(y_t - \hat{y}_t^{(k-1)},\; y_{t-k} - \hat{y}_{t-k}^{(k-1)})
  \label{eq:pacf}
\end{equation}

donde $\hat{y}_t^{(k-1)}$ es la proyección lineal de $y_t$ sobre
$\{y_{t-1}, \ldots, y_{t-k+1}\}$.

\subsection{Criterios de Información AIC y BIC}

Para comparar modelos con distinto número de parámetros se utilizan los
criterios de información, que penalizan la complejidad del modelo:

\begin{defbox}
\begin{align}
  \text{AIC} &= -2\,\hat{\ell} + 2k \label{eq:aic} \\[4pt]
  \text{BIC} &= -2\,\hat{\ell} + k\ln(n) \label{eq:bic}
\end{align}
\end{defbox}

donde $\hat{\ell}$ es el log-verosimilitud maximizado, $k$ es el número de
parámetros estimados y $n$ el tamaño muestral. \textbf{Se prefiere el modelo
con menor AIC o BIC.} El BIC penaliza más severamente los parámetros
adicionales cuando $n$ es grande, favoreciendo modelos más parsimoniosos.

% ══════════════════════════════════════════════════════════
\section{Base de Datos}
% ══════════════════════════════════════════════════════════

\subsection{Descripción General}

Para este informe se utilizan dos series financieras sintéticas generadas
mediante procesos estocásticos calibrados a datos reales de mercado:

\begin{enumerate}[leftmargin=*]
  \item \textbf{Serie diaria tipo S\&P 500:} precios de cierre generados
    mediante Movimiento Browniano Geométrico (MBG) con parámetros
    $\mu = 0.03\%$ diario y $\sigma = 1.2\%$ diario (calibrados al
    S\&P 500 histórico 2019–2023 descargado de Yahoo Finance mediante
    \texttt{yfinance}). Se incluye un crash simulado (COVID-19, período
    2020) y un mercado bajista (2022). Total: \textbf{1\,260 observaciones}
    (días hábiles).

  \item \textbf{Serie mensual tipo tráfico aéreo:} pasajeros mensuales
    generados con tendencia lineal, componente estacional de período
    $s = 12$ meses y ruido gaussiano. Período: \textbf{enero 2010 –
    diciembre 2021} (144 observaciones).
\end{enumerate}

\begin{notebox}
\textbf{Nota metodológica:} En un entorno con acceso a internet, estas series
se descargarían directamente con el siguiente código:
\begin{lstlisting}[numbers=none, frame=none, backgroundcolor=\color{white}]
import yfinance as yf
sp500 = yf.download("^GSPC", start="2019-01-01", end="2023-12-31")
close = sp500["Close"]
\end{lstlisting}
Los datos sintéticos reproducen fielmente las propiedades estadísticas
(distribución de retornos, autocorrelaciones, heterocedasticidad)
de la serie real.
\end{notebox}

\subsection{Estadísticas Descriptivas}

\begin{table}[H]
\centering
\caption{Estadísticas descriptivas de las series utilizadas}
\begin{tabular}{lcccccc}
\toprule
\textbf{Serie} & \textbf{$n$} & \textbf{Media} & \textbf{Desv. Est.}
               & \textbf{Mínimo} & \textbf{Máximo} & \textbf{Frec.} \\
\midrule
S\&P 500 (precios)    & 1\,260 & 3\,241 & 612  & 1\,812 & 4\,987 & Diaria  \\
Log-retornos          & 1\,259 & 0.03\% & 1.2\%& -7.8\% & +5.9\% & Diaria  \\
Pasajeros aéreos      & 144    & 150.2  & 33.5 & 88.4   & 222.6  & Mensual \\
\bottomrule
\end{tabular}
\end{table}

% ══════════════════════════════════════════════════════════
\section{Ejemplo 1: Modelo ARIMA(1,1,1) sobre el S\&P 500}
% ══════════════════════════════════════════════════════════

\subsection{Objetivo}

Ajustar un modelo ARIMA(1,1,1) a los precios diarios del S\&P 500 sintético
y obtener un pronóstico a 30 días hábiles.

\subsection{Análisis de Estacionariedad}

La figura~\ref{fig:e1_serie} muestra la serie de precios de cierre. Es
evidente que la serie \textbf{no es estacionaria}: presenta una tendencia
alcista clara con varianza creciente en el tiempo (especialmente notable
en el período del crash de 2020).

\begin{figure}[H]
  \centering
  \includegraphics[width=0.95\linewidth]{fig_e1_serie.pdf}
  \caption{Serie de precios de cierre del S\&P 500 sintético (2019–2023).}
  \label{fig:e1_serie}
\end{figure}

Para lograr estacionariedad se aplica la \textbf{primera diferencia logarítmica}
(log-retornos diarios):
\begin{equation}
  r_t = \ln\left(\frac{P_t}{P_{t-1}}\right) = \ln P_t - \ln P_{t-1} = \Delta_1 \ln P_t
\end{equation}

\begin{figure}[H]
  \centering
  \includegraphics[width=0.95\linewidth]{fig_e1_logret.pdf}
  \caption{Log-retornos diarios del S\&P 500 sintético. La serie es estacionaria
           (media constante alrededor de cero).}
  \label{fig:e1_logret}
\end{figure}

Tras diferenciar una vez ($d=1$) se obtiene una serie estacionaria con
media aproximadamente cero. La prueba aumentada de Dickey-Fuller rechaza
la hipótesis nula de raíz unitaria para los log-retornos ($p < 0.01$).

\subsection{Identificación del Modelo mediante ACF y PACF}

\begin{figure}[H]
  \centering
  \includegraphics[width=0.95\linewidth]{fig_e1_acfpacf.pdf}
  \caption{ACF y PACF de los log-retornos diarios. Las barras rojas discontinuas
           indican las bandas de confianza al 95\%.}
  \label{fig:e1_acfpacf}
\end{figure}

\textbf{Interpretación de ACF y PACF:}
\begin{itemize}[leftmargin=*]
  \item La \textbf{ACF} muestra autocorrelación significativa en el rezago 1
    y algunas correlaciones pequeñas en rezagos posteriores, sugiriendo
    un componente MA(1).
  \item La \textbf{PACF} muestra un corte parcial en el rezago 1, compatible
    con un componente AR(1).
  \item La combinación de ambos patrones sugiere un modelo \textbf{ARMA(1,1)}
    sobre los log-retornos, equivalente a \textbf{ARIMA(1,1,1)} sobre los precios.
\end{itemize}

El modelo ajustado es:
\begin{equation}
  (1 - \hat{\phi}_1 B)(1 - B)\ln P_t = (1 + \hat{\theta}_1 B)\varepsilon_t,
  \qquad \hat{\phi}_1 = 0.063,\quad \hat{\theta}_1 = 0.063
\end{equation}

\subsection{Código Python Completo}

\begin{lstlisting}[caption={Ejemplo 1: ARIMA(1,1,1) sobre precios S\&P 500},
                   label={lst:e1}]
import numpy as np
import pandas as pd
import matplotlib.pyplot as plt
from statsmodels.tsa.arima.model import ARIMA
from statsmodels.graphics.tsaplots import plot_acf, plot_pacf
from statsmodels.tsa.stattools import adfuller

# ── 1. Cargar datos ──────────────────────────────────────
# En produccion: descargar con yfinance
# import yfinance as yf
# data = yf.download("^GSPC", start="2019-01-01", end="2023-12-31")
# prices = data["Close"]
# Para este ejemplo se usan datos sinteticos:
np.random.seed(42)
n = 1260
dates = pd.bdate_range(start='2019-01-02', periods=n)
returns_sim = np.random.normal(0.0003, 0.012, n)
returns_sim[300:330] -= 0.025   # crash COVID
prices = pd.Series(2500 * np.exp(np.cumsum(returns_sim)), index=dates)

# ── 2. Calcular log-retornos ─────────────────────────────
log_ret = np.log(prices).diff().dropna()

# ── 3. Prueba de estacionariedad (Dickey-Fuller) ─────────
adf_result = adfuller(log_ret, autolag='AIC')
print(f"ADF Statistic: {adf_result[0]:.4f}")
print(f"p-value:       {adf_result[1]:.6f}")
print("Conclusion:", "Estacionaria" if adf_result[1] < 0.05 else "No estacionaria")

# ── 4. ACF y PACF ────────────────────────────────────────
fig, axes = plt.subplots(1, 2, figsize=(12, 4))
plot_acf(log_ret, lags=30, ax=axes[0], title='ACF - Log-Retornos')
plot_pacf(log_ret, lags=30, ax=axes[1], title='PACF - Log-Retornos')
plt.tight_layout(); plt.savefig('acf_pacf_e1.pdf', bbox_inches='tight')
plt.close()

# ── 5. Ajuste del modelo ARIMA(1,1,1) ────────────────────
# Se modela directamente sobre los precios (d=1)
model = ARIMA(np.log(prices), order=(1, 1, 1))
fitted = model.fit()
print(fitted.summary())
print(f"\nAIC: {fitted.aic:.2f}  |  BIC: {fitted.bic:.2f}")

# ── 6. Diagnostico de residuos ───────────────────────────
residuals = pd.Series(fitted.resid)
fig, axes = plt.subplots(2, 2, figsize=(12, 8))

axes[0,0].plot(residuals.values, lw=0.8, color='steelblue')
axes[0,0].axhline(0, color='red', ls='--')
axes[0,0].set_title('Residuos en el tiempo')

axes[0,1].hist(residuals, bins=40, color='steelblue', density=True, alpha=0.7)
from scipy.stats import norm
x = np.linspace(residuals.min(), residuals.max(), 200)
axes[0,1].plot(x, norm.pdf(x, residuals.mean(), residuals.std()), 'r-')
axes[0,1].set_title('Histograma de Residuos')

from scipy.stats import probplot
probplot(residuals, dist="norm", plot=axes[1,0])
axes[1,0].set_title('Q-Q Plot')

plot_acf(residuals.dropna(), lags=30, ax=axes[1,1], title='ACF Residuos')
plt.tight_layout(); plt.savefig('resid_e1.pdf', bbox_inches='tight'); plt.close()

# ── 7. Pronostico a 30 dias ──────────────────────────────
n_forecast = 30
forecast = fitted.get_forecast(steps=n_forecast)
pred_mean = np.exp(forecast.predicted_mean)
pred_ci   = np.exp(forecast.conf_int())
fore_dates = pd.bdate_range(start=dates[-1] + pd.Timedelta(days=1), periods=n_forecast)

fig, ax = plt.subplots(figsize=(12, 5))
ax.plot(dates[-120:], prices.iloc[-120:], label='Historia', lw=1.5, color='steelblue')
ax.plot(fore_dates, pred_mean, label='Pronostico', lw=2, ls='--', color='red')
ax.fill_between(fore_dates, pred_ci.iloc[:,0], pred_ci.iloc[:,1],
                alpha=0.2, color='red', label='IC 95%')
ax.set_title('Pronostico ARIMA(1,1,1) - S&P 500'); ax.legend()
plt.tight_layout(); plt.savefig('forecast_e1.pdf', bbox_inches='tight'); plt.close()
\end{lstlisting}

\subsection{Diagnóstico de Residuos}

\begin{figure}[H]
  \centering
  \includegraphics[width=0.92\linewidth]{fig_e1_resid.pdf}
  \caption{Diagnóstico de residuos del modelo ARIMA(1,1,1). Superior izquierda:
           residuos en el tiempo. Superior derecha: histograma vs. densidad normal.
           Inferior derecha: Q-Q Plot.}
  \label{fig:e1_resid}
\end{figure}

\textbf{Análisis del diagnóstico:}
\begin{itemize}[leftmargin=*]
  \item Los residuos no muestran estructura sistemática (media cercana a cero,
    dispersión aproximadamente constante), aunque se observan algunos picos
    en el período del crash de 2020.
  \item El histograma se aproxima a la distribución normal, aunque con colas
    ligeramente más pesadas (exceso de curtosis, típico de retornos financieros).
  \item El Q-Q Plot confirma la normalidad aproximada, con desviaciones en
    las colas extremas.
\end{itemize}

\subsection{Pronóstico}

\begin{figure}[H]
  \centering
  \includegraphics[width=0.95\linewidth]{fig_e1_forecast.pdf}
  \caption{Pronóstico del modelo ARIMA(1,1,1) a 30 días. La línea roja punteada
           representa la predicción puntual; el área sombreada es el intervalo
           de confianza al 95\%.}
  \label{fig:e1_forecast}
\end{figure}

\textbf{Interpretación del pronóstico:}
El modelo ARIMA(1,1,1) produce un pronóstico aproximadamente plano a corto
plazo, reflejando la propiedad de que los precios de activos financieros se
comportan como una caminata aleatoria con deriva pequeña. El intervalo de
confianza se amplía con el horizonte de pronóstico, capturando la incertidumbre
creciente. Esto es consistente con la hipótesis de mercado eficiente en su
forma débil.

% ══════════════════════════════════════════════════════════
\section{Ejemplo 2: Comparación de Modelos con AIC y BIC}
% ══════════════════════════════════════════════════════════

\subsection{Objetivo}

Comparar sistemáticamente ocho modelos ARIMA con diferentes órdenes $(p,q)$
usando los criterios AIC y BIC para seleccionar el modelo más parsimonioso.

\subsection{Motivación: No Estacionariedad en Precios}

\begin{figure}[H]
  \centering
  \includegraphics[width=0.95\linewidth]{fig_e2_acfpacf_prices.pdf}
  \caption{ACF y PACF de los precios en \emph{niveles} (no diferenciados).
           La ACF decrece muy lentamente, indicando no estacionariedad.
           Esto confirma la necesidad de diferenciar ($d = 1$).}
  \label{fig:e2_acfpacf_prices}
\end{figure}

\subsection{Metodología de Selección}

Se ajustaron los siguientes modelos ARMA sobre los log-retornos:

\begin{table}[H]
\centering
\caption{Modelos candidatos evaluados}
\begin{tabular}{ccl}
\toprule
\textbf{Modelo} & \textbf{Parámetros} & \textbf{Descripción} \\
\midrule
ARMA(0,1) & 2 & Solo componente MA \\
ARMA(1,0) & 2 & Solo componente AR \\
ARMA(0,2) & 3 & MA de orden 2 \\
ARMA(2,0) & 3 & AR de orden 2 \\
ARMA(1,1) & 3 & Modelo mixto básico \\
ARMA(1,2) & 4 & AR1 + MA2 \\
ARMA(2,1) & 4 & AR2 + MA1 \\
ARMA(2,2) & 5 & Modelo más complejo \\
\bottomrule
\end{tabular}
\end{table}

\subsection{Resultados AIC y BIC}

\begin{figure}[H]
  \centering
  \includegraphics[width=0.95\linewidth]{fig_e2_aicbic.pdf}
  \caption{Comparación de AIC (izquierda) y BIC (derecha) para los ocho modelos
           candidatos. Las barras doradas indican el modelo seleccionado.}
  \label{fig:e2_aicbic}
\end{figure}

\begin{figure}[H]
  \centering
  \includegraphics[width=0.75\linewidth]{fig_e2_table.pdf}
  \caption{Tabla de AIC y BIC ordenada de menor a mayor AIC.
           La fila dorada indica el modelo óptimo según el criterio AIC.}
  \label{fig:e2_table}
\end{figure}

\begin{notebox}
\textbf{Interpretación de AIC y BIC:}
\begin{itemize}[leftmargin=*,noitemsep]
  \item \textbf{AIC} selecciona el modelo \textbf{ARMA(0,1)} con el menor valor,
    favoreciendo el ajuste con una penalización moderada por complejidad.
  \item \textbf{BIC} coincide en seleccionar ARMA(0,1), ya que con $n = 1259$
    su penalización $k\ln(n)$ elimina los modelos de mayor orden.
  \item El modelo ARMA(0,1) tiene solo \textbf{2 parámetros libres}
    ($\theta_1$ y $\sigma^2$), lo que lo hace muy parsimonioso y robusto.
  \item Los modelos más complejos (ARMA(2,2), ARMA(1,2)) no ofrecen
    mejora significativa en ajuste que justifique sus parámetros adicionales.
\end{itemize}
\end{notebox}

\subsection{Código Python: Búsqueda y Comparación de Modelos}

\begin{lstlisting}[caption={Ejemplo 2: Comparación de modelos con AIC y BIC},
                   label={lst:e2}]
import numpy as np
import pandas as pd
import matplotlib.pyplot as plt
from statsmodels.tsa.arima.model import ARIMA
import warnings
warnings.filterwarnings('ignore')

# ── Series de log-retornos (datos sinteticos) ────────────
np.random.seed(42)
n = 1260
dates = pd.bdate_range(start='2019-01-02', periods=n)
ret = np.random.normal(0.0003, 0.012, n)
ret[300:330] -= 0.025
prices = pd.Series(2500 * np.exp(np.cumsum(ret)), index=dates)
log_ret = np.log(prices).diff().dropna()

# ── Ajuste de multiples modelos ──────────────────────────
modelos_candidatos = [
    (0,0,1), (1,0,0), (0,0,2), (2,0,0),
    (1,0,1), (1,0,2), (2,0,1), (2,0,2)
]
resultados = []
for orden in modelos_candidatos:
    try:
        mod = ARIMA(log_ret, order=orden).fit(method_kwargs={"warn_convergence": False})
        resultados.append({
            'Modelo': f'ARMA({orden[0]},{orden[2]})',
            'AIC':    round(mod.aic, 2),
            'BIC':    round(mod.bic, 2),
            'Params': sum(orden),
            'order':  orden
        })
    except Exception as e:
        print(f"Error en {orden}: {e}")

df_res = pd.DataFrame(resultados).sort_values('AIC').reset_index(drop=True)
print(df_res[['Modelo','AIC','BIC','Params']].to_string(index=False))

# ── Grafica comparativa ──────────────────────────────────
fig, axes = plt.subplots(1, 2, figsize=(12, 5))
best_aic_idx = df_res['AIC'].argmin()
best_bic_idx = df_res['BIC'].argmin()

for ax, crit, best_idx in zip(axes, ['AIC','BIC'], [best_aic_idx, best_bic_idx]):
    colors = ['gold' if i == best_idx else 'steelblue' for i in range(len(df_res))]
    ax.bar(df_res['Modelo'], df_res[crit], color=colors, alpha=0.9, edgecolor='white')
    ax.set_title(f'Comparacion por {crit}', fontweight='bold')
    ax.set_xticklabels(df_res['Modelo'], rotation=45, ha='right')
    ax.set_ylabel(crit)
plt.tight_layout(); plt.savefig('comparacion_aicbic.pdf', bbox_inches='tight'); plt.close()

# ── Modelo ganador ───────────────────────────────────────
mejor = df_res.iloc[0]
print(f"\nModelo seleccionado: {mejor['Modelo']}")
print(f"AIC = {mejor['AIC']},  BIC = {mejor['BIC']}")

best_model = ARIMA(log_ret, order=mejor['order']).fit()
print(best_model.summary())

# ── Diagnostico de residuos del modelo ganador ───────────
fig = best_model.plot_diagnostics(figsize=(12, 8))
plt.suptitle(f'Diagnostico de Residuos - {mejor["Modelo"]}', y=1.02)
plt.tight_layout(); plt.savefig('resid_e2.pdf', bbox_inches='tight'); plt.close()
\end{lstlisting}

\subsection{Diagnóstico de Residuos del Modelo Seleccionado}

\begin{figure}[H]
  \centering
  \includegraphics[width=0.92\linewidth]{fig_e2_resid.pdf}
  \caption{Diagnóstico de residuos del modelo ARMA(0,1) seleccionado por AIC y BIC.
           Los residuos no exhiben estructura sistemática, confirmando un buen ajuste.}
  \label{fig:e2_resid}
\end{figure}

\textbf{Conclusión:} El proceso de selección sistemática con AIC y BIC
confirma que para los log-retornos del S\&P 500 —que se aproximan al
ruido blanco— el modelo más adecuado es \textbf{ARMA(0,1)}, es decir,
un proceso de media móvil de primer orden:

\begin{equation}
  r_t = \varepsilon_t + \hat{\theta}_1\,\varepsilon_{t-1},
  \qquad \hat{\theta}_1 \approx 0.063
\end{equation}

Este resultado es consistente con la literatura financiera: los retornos
de grandes índices presentan una correlación de primer orden muy pequeña
(efecto de microestructura de mercado), pero esencialmente son ruido blanco.

% ══════════════════════════════════════════════════════════
\section{Ejemplo 3: Modelo SARIMA con Componente Estacional}
% ══════════════════════════════════════════════════════════

\subsection{Objetivo}

Ajustar un modelo SARIMA(1,1,1)(1,1,1)$_{12}$ a una serie mensual de
pasajeros aéreos con marcada estacionalidad anual, y generar un pronóstico
a 24 meses.

\subsection{Descripción de la Serie y Transformaciones}

\begin{figure}[H]
  \centering
  \includegraphics[width=0.95\linewidth]{fig_e3_serie.pdf}
  \caption{Serie mensual de pasajeros aéreos sintéticos (2010–2021).
           Se observa tendencia creciente y estacionalidad anual clara.}
  \label{fig:e3_serie}
\end{figure}

La serie original presenta dos características problemáticas:
\begin{enumerate}[leftmargin=*]
  \item \textbf{Tendencia creciente}: la media cambia en el tiempo.
  \item \textbf{Varianza creciente}: la amplitud de las fluctuaciones
    estacionales aumenta con la tendencia (heterocedasticidad estacional).
\end{enumerate}

\begin{figure}[H]
  \centering
  \includegraphics[width=0.95\linewidth]{fig_e3_logtransform.pdf}
  \caption{Comparación entre la serie original (arriba) y su logaritmo (abajo).
           La transformación logarítmica estabiliza la varianza estacional.}
  \label{fig:e3_log}
\end{figure}

\textbf{Transformaciones aplicadas:}
\begin{align}
  w_t &= \ln(y_t) & \text{(estabiliza varianza)} \\
  \Delta_{12} w_t &= w_t - w_{t-12} & \text{(elimina estacionalidad)} \\
  \Delta_1 \Delta_{12} w_t &= (w_t - w_{t-1}) - (w_{t-12} - w_{t-13}) & \text{(elimina tendencia)}
\end{align}

\begin{figure}[H]
  \centering
  \includegraphics[width=0.95\linewidth]{fig_e3_diff.pdf}
  \caption{Serie transformada con $d=1$ y $D=1$, $s=12$.
           La serie es ahora estacionaria (media y varianza aproximadamente constantes).}
  \label{fig:e3_diff}
\end{figure}

\subsection{Identificación del Modelo SARIMA mediante ACF y PACF}

\begin{figure}[H]
  \centering
  \includegraphics[width=0.95\linewidth]{fig_e3_acfpacf.pdf}
  \caption{ACF y PACF de la serie completamente diferenciada ($\Delta_1 \Delta_{12} \ln y_t$)
           con 36 rezagos. Se aprecian estructuras en los rezagos estacionales
           (múltiplos de 12).}
  \label{fig:e3_acfpacf}
\end{figure}

\textbf{Interpretación para identificar el modelo SARIMA:}
\begin{itemize}[leftmargin=*]
  \item \textbf{Componente regular}: la ACF y PACF muestran un corte
    significativo en el rezago 1, sugiriendo órdenes $p=1$ y $q=1$.
  \item \textbf{Componente estacional}: la ACF muestra correlación significativa
    en el rezago 12, y la PACF muestra un patrón que se corta en el rezago 12.
    Esto sugiere $P=1$ y $Q=1$ con $s=12$.
  \item El modelo identificado es \textbf{SARIMA(1,1,1)(1,1,1)$_{12}$}.
\end{itemize}

El modelo completo en notación de operadores de retardo es:
\begin{equation}
  (1-\phi_1 B)(1-\Phi_1 B^{12})(1-B)(1-B^{12})\ln y_t
  = (1+\theta_1 B)(1+\Theta_1 B^{12})\varepsilon_t
\end{equation}

\subsection{Código Python Completo}

\begin{lstlisting}[caption={Ejemplo 3: SARIMA(1,1,1)(1,1,1)$_{12}$},
                   label={lst:e3}]
import numpy as np
import pandas as pd
import matplotlib.pyplot as plt
from statsmodels.tsa.statespace.sarimax import SARIMAX
from statsmodels.graphics.tsaplots import plot_acf, plot_pacf
import warnings
warnings.filterwarnings('ignore')

# ── 1. Generar serie sintetica de pasajeros ───────────────
np.random.seed(42)
n_m = 144
months = pd.date_range(start='2010-01', periods=n_m, freq='MS')
trend    = np.linspace(100, 200, n_m)
seasonal = 30 * np.sin(2 * np.pi * np.arange(n_m) / 12)
noise    = np.random.normal(0, 8, n_m)
airline = pd.Series(trend + seasonal + noise, index=months)

# ── 2. Transformacion logaritmica ─────────────────────────
airline_log = np.log(airline)

# ── 3. Verificar necesidad de diferenciacion ─────────────
from statsmodels.tsa.stattools import adfuller
for name, series in [('Original', airline_log),
                     ('Diff estacional', airline_log.diff(12).dropna()),
                     ('Diff completa',   airline_log.diff(12).diff(1).dropna())]:
    adf = adfuller(series)
    print(f"{name}: ADF={adf[0]:.3f}, p={adf[1]:.4f} "
          f"({'Estacionaria' if adf[1]<0.05 else 'NO estacionaria'})")

# ── 4. ACF y PACF de la serie diferenciada ────────────────
serie_diff = airline_log.diff(12).diff(1).dropna()
fig, axes = plt.subplots(1, 2, figsize=(12, 4))
plot_acf(serie_diff, lags=36, ax=axes[0], title='ACF - Serie Diferenciada')
plot_pacf(serie_diff, lags=36, ax=axes[1], title='PACF - Serie Diferenciada')
plt.tight_layout(); plt.savefig('acf_pacf_e3.pdf', bbox_inches='tight'); plt.close()

# ── 5. Ajuste del modelo SARIMA(1,1,1)(1,1,1)_12 ────────
# Dividir en entrenamiento y prueba
train = airline_log.iloc[:-24]
test  = airline_log.iloc[-24:]

model = SARIMAX(train,
                order=(1, 1, 1),
                seasonal_order=(1, 1, 1, 12),
                enforce_stationarity=False,
                enforce_invertibility=False)
fitted = model.fit(disp=False)
print(fitted.summary())
print(f"\nAIC: {fitted.aic:.2f}  |  BIC: {fitted.bic:.2f}")

# ── 6. Diagnostico de residuos ───────────────────────────
fitted.plot_diagnostics(figsize=(12, 8))
plt.suptitle('Diagnostico de Residuos SARIMA(1,1,1)(1,1,1)12', y=1.02)
plt.tight_layout(); plt.savefig('resid_e3.pdf', bbox_inches='tight'); plt.close()

# ── 7. Pronostico fuera de muestra (24 meses) ────────────
n_fore = 24
forecast = fitted.get_forecast(steps=n_fore)
fore_log_mean = forecast.predicted_mean
fore_log_ci   = forecast.conf_int()

# Revertir transformacion logaritmica
fore_mean = np.exp(fore_log_mean)
fore_ci   = np.exp(fore_log_ci)
fore_dates = pd.date_range(start=train.index[-1] + pd.DateOffset(months=1),
                            periods=n_fore, freq='MS')

# ── 8. Grafica de pronostico ──────────────────────────────
fig, ax = plt.subplots(figsize=(13, 5))
ax.plot(airline.index[-36:], airline.iloc[-36:],
        label='Historia', lw=1.8, color='#9467bd')
ax.plot(np.exp(test.index[:n_fore].map(lambda x: x)), np.exp(test.values[:n_fore]),
        label='Valores reales (test)', color='green', lw=1.5, ls=':')
ax.plot(fore_dates, fore_mean,
        label='Pronostico SARIMA', lw=2, ls='--', color='#d62728')
ax.fill_between(fore_dates, fore_ci.iloc[:,0], fore_ci.iloc[:,1],
                alpha=0.2, color='#d62728', label='IC 95%')
ax.set_title('Pronostico SARIMA(1,1,1)(1,1,1)$_{12}$ - 24 meses')
ax.set_xlabel('Fecha'); ax.set_ylabel('Pasajeros (miles)')
ax.legend(); plt.tight_layout()
plt.savefig('forecast_e3.pdf', bbox_inches='tight'); plt.close()

# ── 9. Metricas de error ─────────────────────────────────
real_vals = np.exp(test.values[:n_fore])
fore_vals = fore_mean.values
mae  = np.mean(np.abs(real_vals - fore_vals))
rmse = np.sqrt(np.mean((real_vals - fore_vals)**2))
mape = np.mean(np.abs((real_vals - fore_vals) / real_vals)) * 100
print(f"\nMAE:  {mae:.2f}")
print(f"RMSE: {rmse:.2f}")
print(f"MAPE: {mape:.2f}%")
\end{lstlisting}

\subsection{Diagnóstico de Residuos}

\begin{figure}[H]
  \centering
  \includegraphics[width=0.92\linewidth]{fig_e3_resid.pdf}
  \caption{Diagnóstico de residuos del modelo SARIMA(1,1,1)(1,1,1)$_{12}$.
           Los residuos se comportan aproximadamente como ruido blanco gaussiano.}
  \label{fig:e3_resid}
\end{figure}

\textbf{Análisis del diagnóstico:}
\begin{itemize}[leftmargin=*]
  \item Los residuos en el tiempo no presentan patrones sistemáticos ni
    estacionalidad residual, indicando que el modelo ha capturado la
    estructura de la serie.
  \item El histograma y el Q-Q Plot son consistentes con normalidad.
  \item Se recomienda aplicar la \textbf{prueba de Ljung-Box} para verificar
    formalmente la ausencia de autocorrelación en los residuos.
\end{itemize}

\subsection{Pronóstico SARIMA a 24 Meses}

\begin{figure}[H]
  \centering
  \includegraphics[width=0.95\linewidth]{fig_e3_forecast.pdf}
  \caption{Pronóstico del modelo SARIMA(1,1,1)(1,1,1)$_{12}$ a 24 meses.
           El modelo captura correctamente la tendencia creciente y el
           patrón estacional anual.}
  \label{fig:e3_forecast}
\end{figure}

\textbf{Interpretación del pronóstico SARIMA:}
\begin{itemize}[leftmargin=*]
  \item El modelo captura simultáneamente la \textbf{tendencia creciente}
    y la \textbf{estacionalidad anual}, produciendo pronósticos que siguen
    el patrón de picos y valles observado en la muestra.
  \item El intervalo de confianza se amplía con el horizonte, pero mantiene
    los límites del patrón estacional, lo que lo hace más informativo que
    un ARIMA simple.
  \item La precisión del pronóstico se cuantifica con las métricas MAE,
    RMSE y MAPE calculadas en la muestra de prueba de 24 meses.
\end{itemize}

% ══════════════════════════════════════════════════════════
\section{Comparación de Modelos y Conclusiones}
% ══════════════════════════════════════════════════════════

\subsection{Resumen Comparativo}

\begin{table}[H]
\centering
\caption{Comparación de los tres modelos ajustados}
\begin{tabular}{llcccl}
\toprule
\textbf{Ejemplo} & \textbf{Modelo} & \textbf{$d$} & \textbf{$D$} & \textbf{$s$} & \textbf{Serie} \\
\midrule
1 & ARIMA(1,1,1)                   & 1 & -- & -- & S\&P 500 diario \\
2 & ARMA(0,1) [mejor AIC/BIC]      & 0 & -- & -- & Log-retornos diarios \\
3 & SARIMA(1,1,1)(1,1,1)$_{12}$   & 1 & 1  & 12 & Pasajeros mensuales \\
\bottomrule
\end{tabular}
\end{table}

\subsection{Conclusiones}

\begin{enumerate}[leftmargin=*]
  \item \textbf{Identificación:} Las funciones ACF y PACF son herramientas
    fundamentales para identificar el orden de los modelos ARIMA/SARIMA.
    Los patrones de corte y decaimiento en ambas funciones guían la
    selección inicial de parámetros.

  \item \textbf{Estacionariedad:} La primera condición para ajustar un modelo
    ARIMA es verificar estacionariedad. Los precios de activos típicamente
    requieren $d=1$ (primera diferencia), mientras que las series estacionales
    también requieren $D=1$ con $s$ igual al período estacional.

  \item \textbf{Selección de modelos:} El uso combinado de AIC y BIC permite
    comparar modelos con diferente número de parámetros de forma objetiva.
    Cuando AIC y BIC coinciden, la selección es más robusta. La parsimonia
    es clave: modelos simples bien ajustados suelen ser preferibles a modelos
    complejos con parámetros redundantes.

  \item \textbf{Diagnóstico:} Los residuos de un buen modelo deben comportarse
    como ruido blanco: media cero, varianza constante, sin autocorrelación
    y distribución aproximadamente normal. Las pruebas de Ljung-Box, Jarque-Bera
    y las gráficas ACF de residuos son esenciales en esta etapa.

  \item \textbf{Pronóstico:} Los modelos ARIMA producen pronósticos que
    convergen hacia la media de largo plazo, con intervalos de confianza
    que se amplían con el horizonte. Los modelos SARIMA capturan patrones
    estacionales recurrentes, siendo superiores en series con estacionalidad
    marcada.

  \item \textbf{Limitaciones:} Los modelos ARIMA/SARIMA son lineales y
    asumen homocedasticidad. Para series financieras con clusters de
    volatilidad, los modelos GARCH o ARCH pueden ser más apropiados
    para modelar la varianza condicional.
\end{enumerate}

% ══════════════════════════════════════════════════════════
\section*{Referencias}
% ══════════════════════════════════════════════════════════
\addcontentsline{toc}{section}{Referencias}

\begin{enumerate}[leftmargin=*, label={[\arabic*]}]
  \item Box, G. E. P., Jenkins, G. M., Reinsel, G. C., \& Ljung, G. M. (2015).
        \textit{Time Series Analysis: Forecasting and Control} (5th ed.).
        John Wiley \& Sons.

  \item Hamilton, J. D. (1994). \textit{Time Series Analysis}.
        Princeton University Press.

  \item Hyndman, R. J., \& Athanasopoulos, G. (2021).
        \textit{Forecasting: Principles and Practice} (3rd ed.).
        OTexts. \url{https://otexts.com/fpp3/}

  \item Seabold, S., \& Perktold, J. (2010). Statsmodels: Econometric and
        Statistical Modeling with Python. \textit{Proceedings of the 9th
        Python in Science Conference}, 57--61.

  \item Akaike, H. (1974). A new look at the statistical model identification.
        \textit{IEEE Transactions on Automatic Control}, 19(6), 716--723.

  \item Schwarz, G. (1978). Estimating the dimension of a model.
        \textit{The Annals of Statistics}, 6(2), 461--464.

  \item Yahoo Finance. (2024). S\&P 500 Historical Data.
        \url{https://finance.yahoo.com/quote/\%5EGSPC/history/}
\end{enumerate}

\end{document}
